\documentclass[instructions.tex]{subfiles}
\begin{document}
 
\chapter{Du prix de notre âme.}
 
\primo On doit conclure de tout ce qui a été dit, que notre âme est plus précieuse que toutes les richesses de la terre, et que la perte d’une âme est un plus grand mal que la destruction de toutes les créatures visibles. Jugeons de la dignité d’une âme par celle des anges. Un ange est si parfait, que tout ce que nous voyons sur la terre et dans les cieux, est moindre en comparaison d’un ange, qu’un grain de poussière en comparaison du soleil. Or, notre âme est presque égale aux anges\footnote{Minuisti eum paulò minùs ab angelis. Ps. 8. 6.}. Et dans le ciel elle sera comme divinisée, semblable à Dieu\footnote{Similes ei erimus. 1. Joan. 3. 2.}.
 
Quelque parfaits que soient les anges, notre âme est si chère à Dieu, qu’il a ordonné aux anges d’en prendre soin, et qu’il n’y en a pas qui n’ait un prince du ciel pour sa garde\footnote{(3) Angelis suis mandavit de te ut custodiant te. Ps. 11. 90.} (3).
 
Le démon, pour tenter le Sauveur, lui offrit tous les royaumes du monde; ce qui nous marque que le démon est si jaloux de la beauté et du prix d’une âme, que, s’il était en son pouvoir, il donnerait tous les empires de l’univers pour en avoir une seule. N’est-il pas honteux, dit un Père, que nous estimions moins notre âme que le démon lui-même ne l’estime?
 
\secundo Notre âme est d’un si grand prix, que le Tout-Puissant n’a pas cru faire trop d’envoyer son Fils sur la terre pour la sauver. Quand il n’y aurait eu qu’un seul homme au monde, son âme est si précieuse, dit saint Chrysostome, qu’il n’aurait pas été indigne de Dieu de s’incarner et de donner sa vie pour elle. Fût-elle l’âme du dernier des hommes, elle est si chère à son Créateur, que, si cet homme craint Dieu, s’il observe sa loi, Dieu anéantirait plutôt les cieux que de la laisser périr, parce qu’ils sont peu de chose en comparaison d’une âme.
 
Les cieux n’ont coûté à Dieu qu’une parole, mais notre âme a coûté à Dieu son Fils; elle a coûté au Fils de Dieu son sang et sa vie; elle vaut donc, en un sens, autant que Dieu, puisqu’elle vaut autant qu’elle a coûté. O corps! s’écrie saint Bernard, que vous êtes honoré de loger une âme, de posséder un hôte si digne! Rendez-lui tout l’honneur qu’il mérite.
 
Si vous aviez recueilli dans un vase le sang du Sauveur, lorsqu’il expirait sur la croix, avec quelle respectueuse attention n’auriez-vous pas conservé ce sang adorable! Devez-vous avoir moins de soin pour conserver votre âme, que pour conserver le sang de Jésus-Christ, puisque Jésus-Christ a plus estimé votre âme que son propre sang? Depuis que j’ai connu que mon âme a été rachetée par le sang d’un Dieu, disait saint Augustin, je suis résolu de la conserver, de ne jamais la vendre au démon par le péché.
 
Votre âme ne vous appartient pas; elle est plus à Dieu qu’à vous. C’est un dépôt dont il vous a chargé de lui rendre compte; conservez-la comme le prix du sang de Jésus-Christ, conservez-la pour vous-même; si vous la perdez, tout sera perdu pour vous\footnote{(1) Quam dabit homo commutationem pro animâ suâ?} (1).
 
\section{Résolutions}
 
\primo Oui, il faut que je sauve mon âme, à quelque prix que ce soit.

\secundo Pour m’y encourager, je me rappellerai de temps en temps qu’elle est l’image de Dieu; qu’elle approche des anges par sa beauté; qu’elle a un prince du ciel pour gardien; que Jésus-Christ a répandu son sang pour la racheter, et que j’en rendrai compte à Dieu, à qui elle appartient. 

\prayer{Mon Dieu, faites-moi connoître tout le prix de mon âme, et soutenez-moi dans la résolution que j’ai prise de la sauver pour l’éternité.}
 
\end{document}
